\documentclass{article}

% NeurIPS 2026 style
\usepackage[final]{neurips_2025}  % or neurips_2026 when available

\usepackage{amsmath,amssymb,amsthm}
\usepackage{natbib}
\usepackage{hyperref}
\usepackage{cleveref}
\usepackage{booktabs}
\usepackage{graphicx}
\usepackage{xcolor}
\usepackage{subcaption}
\usepackage{algorithm}
\usepackage{algorithmic}
\usepackage{multirow}

\newcommand{\R}{\mathbb{R}}
\newcommand{\E}{\mathbb{E}}
\newcommand{\bx}{\mathbf{x}}
\newcommand{\ba}{\mathbf{a}}
\newcommand{\bv}{\mathbf{v}}
\newcommand{\bh}{\mathbf{h}}

\title{Stability-Guided Activation Steering in Diffusion Transformers}

\author{
  Author 1 \And Author 2
}

\begin{document}
\maketitle

\begin{abstract}
We study activation steering in Diffusion Transformers (DiTs)---adding learned concept vectors to intermediate activations during denoising to control generation at zero inference cost.
We develop a \emph{stability-theoretic framework} explaining which extraction methods and layer configurations yield effective steering.
Our key findings:
(1)~mean-difference vectors outperform sophisticated alternatives (RFM, logistic regression) because they maximize \emph{temporal stability}---consistency of the concept direction across denoising steps ($S = 0.86$ vs.\ $0.50$);
(2)~multi-layer steering accumulates according to a front-weighted formula validated on 25 layer configurations ($r = 0.95$), reflecting the interaction of perturbation amplification and norm clipping;
(3)~Recursive Feature Machine (AGOP) directions are orthogonal to the causal steering direction ($\cos \approx 0$), explaining their failure in DiTs despite success in other architectures;
(4)~norm-scaling, used by concurrent work, collapses sample diversity to $< 10\%$ of baseline while achieving comparable concept shift.
On DiT-XL/2-256, our method achieves $+0.12$--$0.27$ CLIP concept lift for semantic concepts (animal, natural) and $+0.58$ pixel lift for statistical concepts (brightness) at FID~$66$, compared to FID~$269$ for norm-scaling approaches.
\end{abstract}

% ==================================================================
\section{Introduction}
\label{sec:intro}

Controlling the attributes of diffusion model outputs---steering toward brighter images, more natural scenes, or specific visual styles---is a fundamental capability for deployment.
Current approaches require either retraining (RL fine-tuning, LoRA adapters) or multiplicative inference cost (particle-based methods like Feynman-Kac steering, classifier guidance).
A zero-cost alternative exists: \emph{activation steering}, where learned concept vectors are added to the model's intermediate representations during generation.
This approach has shown remarkable success in large language models~\citep{turner2023activation,zou2023representation,marks2024geometry}, but lacks theoretical grounding in diffusion models.

Concurrent work by \citet{wang2026general} demonstrates that Recursive Feature Machine (RFM) directions can steer unconditional U-Net diffusion models, achieving 96.6\% class accuracy on CIFAR-10.
However, they provide no theoretical explanation for their results and note that ``it would be interesting to see why discriminative directions are effective for generation.''
We provide exactly this explanation---and show that their approach fails on Diffusion Transformers (DiTs), the dominant architecture for modern diffusion models.

\paragraph{Contributions.}
\begin{enumerate}[leftmargin=*,nosep]
  \item \textbf{Temporal stability as the key predictor.} We show that the effectiveness of a steering vector depends primarily on its \emph{temporal stability}---the consistency of its direction across denoising timesteps. Mean-difference vectors have 71\% higher stability than logistic regression vectors ($S = 0.86$ vs.\ $0.50$ at layer~0), explaining their superior causal effectiveness despite lower classification accuracy.

  \item \textbf{Front-weighted accumulation.} We derive and validate a formula predicting multi-layer steering effectiveness from per-layer properties, achieving $r = 0.95$ correlation on 25 held-out layer configurations. The key insight: perturbations amplify through layers (measured Jacobian $\alpha > 1$), but norm clipping bounds the effective contribution, creating a front-weighted pattern.

  \item \textbf{RFM orthogonality in DiTs.} We discover that AGOP/RFM directions are \emph{orthogonal} to mean-difference directions in DiT ($\cos \approx 0$ at all layers), explaining why RFM gives negative steering lift in DiTs despite positive results in U-Nets. This architecture-dependent failure has not been previously documented.

  \item \textbf{Norm-scaling destroys diversity.} We show that the norm-scaling formula $\bh' = \bh + w \|\bh\| \bv$ used by \citet{wang2026general} collapses sample diversity to $<10\%$ of baseline (FID degrades from 83 to 269) while our norm-clipping approach preserves diversity at 93\%.
\end{enumerate}

% ==================================================================
\section{Related Work}
\label{sec:related}

\paragraph{Activation steering in LLMs.}
\citet{turner2023activation} demonstrated that adding ``steering vectors'' to LLM residual streams controls model behavior.
\citet{marks2024geometry} showed that mean-difference directions are more causally effective than logistic regression probes, establishing the predictive-causal gap.
\citet{zou2023representation} formalized representation engineering as a general framework.

\paragraph{Activation editing in diffusion models.}
\citet{kwon2023diffusion} identified h-space (U-Net bottleneck activations) as a semantic latent space amenable to linear editing.
\citet{wang2026general} proposed NA-RFM, using AGOP-based concept vectors for unconditional U-Net diffusion, achieving strong class-conditional generation.
Apple's Activation Transport (AcT)~\citep{kramar2025controlling} demonstrated cross-domain transfer of steering between LLMs and diffusion models.

\paragraph{Inference-time guidance.}
Particle-based methods including Feynman-Kac steering~\citep{phillips2025general} and DriftLite~\citep{ren2026driftlite} provide principled but expensive inference-time control.
Classifier-free guidance (CFG)~\citep{ho2022classifier} is the standard approach but cannot target arbitrary concepts.

\paragraph{Diffusion Transformers.}
DiT~\citep{peebles2023scalable} replaces the U-Net backbone with a vision transformer, using adaptive Layer Normalization (adaLN-Zero) for class conditioning.
DiTs power modern systems including Stable Diffusion 3 and Flux.

% ==================================================================
\section{Method}
\label{sec:method}

\subsection{Concept Vector Extraction}

Given a set of $M$ generated images with activations $\{\ba_i^{(\ell)}\}_{i=1}^M$ at layer $\ell$ and binary concept labels $\{y_i\}$, we extract concept directions using four methods:

\textbf{Mean-difference:} $\bv_\ell = \bar{\ba}_+^{(\ell)} - \bar{\ba}_-^{(\ell)}$, normalized to unit length, where $\bar{\ba}_\pm$ are class-conditional means.

\textbf{PCA:} Top principal component of positive-class activations (centered).

\textbf{Logistic regression:} Weight vector $\mathbf{w}$ from $\arg\max_\mathbf{w} \sum_i \log \sigma(y_i \mathbf{w}^\top \ba_i)$.

\textbf{RFM/AGOP:} Top eigenvector of the Average Gradient Outer Product after $T$ iterations of kernel ridge regression with Laplace kernel~\citep{radhakrishnan2024mechanism}.

\subsection{Steering Mechanism}

At each denoising step, we modify the hidden state after each transformer block $\ell$:
\begin{equation}
  \bh_\ell' = \text{clip}\big(\bh_\ell + \varepsilon \cdot \bv_\ell,\; C \cdot \|\bh_\ell\|\big)
  \label{eq:steer}
\end{equation}
where $\varepsilon$ is the steering strength, $\bv_\ell$ is the unit concept vector at layer $\ell$, and $\text{clip}(\cdot, r)$ rescales the output to have norm at most $r$.
We apply this at \emph{all} $L = 28$ layers simultaneously (``all-layer steering''), using a fixed concept vector shared across all denoising steps (Strategy~A).

The norm clipping parameter $C$ is critical: it bounds the effective perturbation magnitude after amplification through subsequent layers, preventing out-of-distribution activations while preserving the concept direction.

\subsection{Contrast with Norm-Scaling}

\citet{wang2026general} use norm-scaling: $\bh' = \bh + w \|\bh\| \bv$, which scales the perturbation by the activation norm without bounding the result.
We show (Section~\ref{sec:wang}) that this destroys diversity and degrades FID by $3\times$.

% ==================================================================
\section{Theoretical Framework}
\label{sec:theory}

\subsection{Temporal Stability}

\begin{definition}[Temporal Stability]
For concept vector $\bv$ at layer $\ell$, the temporal stability is the mean pairwise cosine similarity of the concept direction across denoising timesteps:
$S(\bv, \ell) = \frac{2}{T(T-1)} \sum_{t_1 < t_2} |\cos(\bv_\ell^{(t_1)}, \bv_\ell^{(t_2)})|$
where $\bv_\ell^{(t)}$ is the concept vector extracted from activations at denoising step $t$.
\end{definition}

A shared (Strategy~A) concept vector is effective only when $S$ is high: if the concept direction varies substantially across timesteps, a fixed perturbation will be aligned with the concept at some steps and misaligned at others, causing partial cancellation.

\textbf{Empirical finding:} Mean-difference vectors have the highest temporal stability at early layers ($S = 0.86$ at layer~0), followed by PCA ($0.59$) and logistic regression ($0.50$).
This directly explains their superior steering effectiveness.

\subsection{Front-Weighted Accumulation}

When steering at a subset of layers $\mathcal{L} \subseteq \{0, \ldots, L-1\}$, the total concept lift follows:
\begin{equation}
  \text{Lift}(\mathcal{L}) \propto \sum_{\ell \in \mathcal{L}} w(\ell), \qquad w(\ell) = \frac{1}{\ell + 1}
  \label{eq:accumulation}
\end{equation}
This front-weighted formula achieves $r = 0.95$ correlation against measured lifts on 25 held-out layer configurations (Figure~\ref{fig:accumulation}).

\textbf{Mechanism.} Perturbations at layer $\ell$ propagate through subsequent layers with measured amplification factors $\alpha_\ell > 1$ (Table~\ref{tab:jacobian}).
Without clipping, early-layer perturbations would grow $6\times$--$46\times$ by the output.
Norm clipping bounds this growth, but the clipping operation reduces the effective magnitude of heavily amplified perturbations.
The net effect: early layers contribute most because their perturbations are amplified to the clipping threshold (full utilization) while maintaining high temporal stability (concept direction preserved).
Late layers contribute less because lower stability causes the post-clipping direction to diverge from the concept.

\subsection{RFM Orthogonality}

\begin{finding}
In DiT-XL/2-256, RFM/AGOP directions are orthogonal to mean-difference directions at all layers: $\cos(\bv_\text{RFM}, \bv_\text{MD}) \approx 0$.
\end{finding}

This explains why RFM gives \emph{negative} steering lift in DiTs: the AGOP-derived direction lies in a subspace orthogonal to the causal steering direction.
The AGOP captures the direction of maximum kernel gradient variation, which in the adaLN-conditioned DiT residual stream does not align with the population centroid displacement.

This is an architecture-dependent phenomenon: in U-Nets, where \citet{wang2026general} report success, the residual stream geometry may create better alignment between AGOP and centroid directions.

% ==================================================================
\section{Experiments}
\label{sec:exp}

\subsection{Setup}

\textbf{Model.} DiT-XL/2-256~\citep{peebles2023scalable} (675M parameters, class-conditional ImageNet $256 \times 256$).
We also present results on Stable Diffusion 1.5 (U-Net) in the appendix.

\textbf{Concepts.} Statistical: brightness (mean pixel $> 0.5$), colorfulness (cross-channel std), warmth (R$-$B channel balance).
Semantic: animal, natural (measured via CLIP zero-shot classification).

\textbf{Metrics.} Pixel lift (change in concept-positive rate), CLIP concept lift (change in CLIP zero-shot score), FID (Inception v3 features vs.\ unsteered reference), diversity ratio (pairwise L2 distance ratio).

\textbf{Extraction.} 500 images, mean-difference vectors at all 28 layers.
50 images suffice for $>75\%$ of maximum lift (data efficiency analysis in appendix).

\subsection{Main Results: Comparison with Wang et al.}
\label{sec:wang}

Table~\ref{tab:wang} compares our approach with \citet{wang2026general}'s components on DiT-XL/2.

\begin{table}[t]
\centering
\caption{Comparison with Wang et al.\ (2026) on DiT-XL/2-256 (brightness concept, 200 images). Our all-layer mean-difference approach achieves $13\times$ higher pixel lift than their RFM, with $3\times$ better FID and preserved diversity.}
\label{tab:wang}
\small
\begin{tabular}{lcccc}
\toprule
Configuration & Pixel Lift & CLIP Lift & FID $\downarrow$ & Diversity \\
\midrule
\textbf{Ours: all-layer MD} & \textbf{+0.640} & $-0.032$ & 92.4 & \textbf{0.93} \\
\textbf{Ours: first-5 MD} & +0.420 & +0.013 & \textbf{83.1} & \textbf{1.05} \\
Ours: all-layer $\varepsilon$=0.1 & +0.235 & $-0.008$ & 88.2 & 1.05 \\
\midrule
Wang: RFM, L25 & +0.050 & +0.002 & 82.8 & 1.00 \\
Wang: MD + norm-scale, L25 & +0.680 & $-0.420$ & 269.0 & 0.20 \\
Wang: MD + norm-scale, L0 & +0.680 & $-0.567$ & 286.6 & 0.08 \\
\bottomrule
\end{tabular}
\end{table}

\subsection{Pareto Frontier}

Figure~\ref{fig:pareto} shows the concept lift vs.\ FID tradeoff at varying $\varepsilon$.

% Placeholder for figure
% \begin{figure}[t]
% \centering
% \includegraphics[width=\linewidth]{figures/pareto.pdf}
% \caption{Pareto frontier: concept lift vs.\ FID. Sweet spot at $\varepsilon = 0.3$--$0.5$ for statistical concepts, $\varepsilon = 2$--$5$ for semantic concepts.}
% \label{fig:pareto}
% \end{figure}

\subsection{Accumulation Validation}

We test 25 non-contiguous layer combinations and compare four accumulation models:

\begin{table}[t]
\centering
\caption{Model comparison for predicting multi-layer steering lift from layer configuration (25 test configurations).}
\label{tab:models}
\small
\begin{tabular}{lc}
\toprule
Model & Correlation ($r$) \\
\midrule
\textbf{Front-weighted: $\sum 1/(\ell+1)$} & \textbf{0.949} \\
Stability $\times$ $\alpha^\ell$: $\sum S(\ell) \cdot 0.773^\ell$ & 0.933 \\
Pure stability: $\sum S(\ell)$ & 0.617 \\
Layer count & 0.538 \\
\bottomrule
\end{tabular}
\end{table}

\subsection{Why Mean-Difference Wins}

\begin{table}[t]
\centering
\caption{Temporal stability and directional alignment explain the effectiveness hierarchy. Measurements at layer~0.}
\label{tab:methods}
\small
\begin{tabular}{lcccc}
\toprule
Method & Stability $S$ & $\cos(\cdot, \text{MD})$ & Pixel Lift & Status \\
\midrule
\textbf{Mean-diff} & \textbf{0.859} & 1.000 & \textbf{+0.525} & Best \\
Logreg & 0.503 & 0.422 & +0.460 & Good \\
PCA & 0.591 & $-0.840$ & $-0.160$ & Inverted \\
RFM & --- & 0.001 & $-0.050$ & Orthogonal \\
\bottomrule
\end{tabular}
\end{table}

\subsection{Semantic Concept Steering}

Using CLIP zero-shot evaluation, we confirm that activation steering controls semantic properties (Table~\ref{tab:semantic}).

\begin{table}[t]
\centering
\caption{CLIP-measured semantic steering (500 images, bootstrap CIs from 3 seeds $\times$ 200 images).}
\label{tab:semantic}
\small
\begin{tabular}{lccc}
\toprule
Concept & $\varepsilon$ & CLIP Lift & FID \\
\midrule
Animal (MD, all-layer) & $+5.0$ & $+0.124 \pm 0.006$ & 144.8 \\
Natural (MD, all-layer) & $-5.0$ & $+0.128 \pm 0.004$ & 159.4 \\
Animal (MD, all-layer) & $+2.0$ & $+0.044$ & 69.2 \\
Natural (MD, all-layer) & $-2.0$ & $+0.031$ & 69.0 \\
\bottomrule
\end{tabular}
\end{table}

% ==================================================================
\section{Discussion}
\label{sec:discussion}

\paragraph{Limitations.}
Our experiments focus on DiT-XL/2-256 with ImageNet class-conditional generation.
The FID degradation at high $\varepsilon$ ($>100$ for semantic concepts) suggests a quality-concept tradeoff that requires further investigation.
CLIP-based brightness evaluation contradicts pixel-based metrics, indicating that mean pixel value is not a perceptually meaningful brightness measure.

\paragraph{Why does RFM fail in DiTs?}
The RFM/AGOP direction captures maximum kernel gradient variation in a Mahalanobis-reweighted feature space.
In DiT's adaLN-conditioned residual stream, this space does not align with the concept centroid displacement direction.
In U-Nets, where conditioning enters through cross-attention and skip connections, the geometry may be more favorable for AGOP alignment.
A systematic study across architectures is needed.

\paragraph{Broader impact.}
Zero-cost activation steering enables controllable generation without retraining, reducing the computational cost of content control.
However, the same mechanism could be used for harmful steering (e.g., generating unsafe content).
The norm-clipping mechanism provides a natural defense: it bounds the magnitude of any perturbation, limiting the degree of controllable deviation from the base model's distribution.

% ==================================================================
\section{Conclusion}

We presented a stability-theoretic framework for understanding activation steering in Diffusion Transformers.
The core insight is simple: mean-difference vectors succeed because they are temporally stable---consistent across denoising timesteps---while discriminative alternatives (RFM, logistic regression) are either orthogonal to the causal direction or too variable to accumulate coherently.
The front-weighted accumulation formula ($r = 0.95$) provides a practical guide for layer selection, and our comparison with concurrent work reveals critical failure modes (diversity collapse from norm-scaling, RFM orthogonality in DiTs) that inform future method design.

\bibliography{references}
\bibliographystyle{plainnat}

% ==================================================================
\appendix
\section{Additional Results}

% TODO: Full ablation tables, NFA analysis, Jacobian measurements,
% U-Net results, multi-step probing, qualitative images, etc.

\end{document}
